\clearpage
\section{Anexo 8.G: Análisis de modos de falla y efectos (FMEA)}
FMEA aplicado a los 15 riesgos con E0 mayor o igual a 0,20. El modo de falla es el evento de T-16 y su efecto es la consecuencia allí declarada. Se utilizan cinco niveles definidos para este proyecto: severidad S=I global (1 corrección localizada, 2 ajuste de paquete, 3 reproceso sin afectar hitos, 4 amenaza a aceptación/servicio/hito, 5 consecuencia sanitaria grave o trazabilidad irrecuperable). O usa el nivel 1--5 de la misma banda de p0/p1, dentro de la ventana del evento, no una frecuencia diaria. D es dificultad para detectar antes del efecto: 1 comprobación automática integral bloqueante; 2 control documental o ensayo completo antes de liberar; 3 muestreo/reconciliación con cobertura parcial; 4 evidencia tardía o validación dispersa; 5 dato histórico no reconstruible o defecto oculto hasta el efecto.

NPR=S por O por D prioriza modos dentro de esta escala (1--125); no equivale a E, a probabilidad ni a HH. Se actúa por severidad 5 aunque NPR sea bajo. La propuesta obliga a tratar S=5, D mayor o igual a 4 y cualquier E alta/crítica. CAL conserva expediente de modo, causa, efecto, prueba y corrección. Las puntuaciones posteriores son proyecciones sujetas a la verificación de la evidencia de cada control. No se usa menor NPR para dispensar requisitos o aceptar un residual alto.

\begingroup
\renewcommand{\thetable}{8.G.1}
\begin{landscape}
\setlength{\textwidth}{\linewidth}
\renewcommand{\arraystretch}{1.08}
\begin{tablalafrox}{FMEA de 15 riesgos altos/críticos y evaluación posterior proyectada. Fuente: eventos T-16 y juicios de planificación de S8.}{tab:cierre-fmea}{G{0.06} F{0.12} F{0.22} F{0.23} F{0.18} Y}{\cab{ID} & \cab{FMEA inicial} & \cab{Modo, disparador y efecto} & \cab{Detección inicial y control} & \cab{FMEA posterior proyectada} & \cab{Acción, responsable y condición de habilitación}}
R-001 & S=4; O=3; D=4; NPR=48. & Cadena del alcance sin certificación antes del límite de 28 días. Dado que la habilitación EDI de cada cadena depende de su participación y certificación externa antes de las cuatro semanas finales de evidencia de E2, entonces puede que una cadena incluida en el alcance no complete su participación o certificación EDI antes de la fecha requerida, teniendo como consecuencia que no se acredite completo el intercambio electrónico exigido en E2 y deban reprogramarse pruebas y coordinación. & D0=4: la validación distribuida puede llegar después de la fecha necesaria. Control/acción: Agenda EDI y certificación de cada cadena, con pruebas de intercambio y registro de pendientes. & S1=4; O1=2; D1=2; NPR1=16. D1=2: acta o ensayo específico previo a habilitación; se invalida la mejora de detección si no existe evidencia. & Al D-5 sin certificación, JP escala en un día hábil y conserva pendiente la cadena; sólo ampliar tras su certificación. Dueño: ARQ; verifica CAL. \\
R-002 & S=5; O=4; D=5; NPR=100. & Lotes o vínculos históricos sin evidencia de origen al perfilar la carga. Dado que la migración sanitaria utiliza antecedentes del ERP, WMS y planillas y requiere vínculos históricos verificables con evidencia de origen, entonces puede que los datos recibidos contengan faltantes de lote o vínculos históricos que no puedan resolverse con esa evidencia, teniendo como consecuencia que la carga sanitaria afectada quede pendiente de conciliación y decisión documentada, y aumente el trabajo de saneamiento. & D0=5: faltante de origen o interfaz no verificable puede permanecer oculto hasta conciliar/aceptar. Control/acción: Inventario de vínculos sanitarios por lote, evidencia de origen y resolución firmada de Calidad sobre cada excepción. & S1=5; O1=4; D1=3; NPR1=60. D1=3: aun conciliando, la cobertura histórica puede ser parcial; no se reconstruye lo inexistente. & Plan adicional: DAT aísla el conjunto el mismo día; Calidad decide antes del corte. No cargar ni despachar lotes sin trazabilidad exigida; si no es reconstruible, elevar decisión de alcance contractual, sin inventar historia. Dueño: DAT; verifica CAL. \\
R-004 & S=4; O=4; D=4; NPR=64. & Diferencias sin causa ni ajuste aprobado en el acta previa al corte. Dado que la aceptación del inventario migrado requiere conciliar los saldos legados con el conteo físico y aprobar las explicaciones y ajustes antes del corte, entonces puede que se detecten diferencias sin explicación y ajuste aprobados a tiempo, teniendo como consecuencia que deban repetirse el saneamiento y la validación, y la carga afectada no pueda ser aceptada hasta resolverlas. & D0=4: la validación distribuida puede llegar después de la fecha necesaria. Control/acción: Acta de conteo, causas de diferencias y ajustes autorizados por Bodega/Finanzas. & S1=4; O1=2; D1=2; NPR1=16. D1=2: acta o ensayo específico previo a habilitación; se invalida la mejora de detección si no existe evidencia. & Bloquear corte con diferencia no explicada; DAT presenta conciliación y nueva prueba antes de la ventana autorizada. Dueño: DAT; verifica CAL. \\
R-005 & S=4; O=3; D=4; NPR=48. & No existe saldo inicial aceptado por cliente y transportista. Dado que la carga inicial de M8 requiere saldos de envases por cliente y transportista sustentados en antecedentes conciliables del cliente, entonces puede que los antecedentes entregados no permitan determinar un saldo conciliado por cliente y transportista, teniendo como consecuencia que deban repetirse la preparación y validación de la carga y no pueda demostrarse M8 con saldos iniciales aceptados. & D0=4: la validación distribuida puede llegar después de la fecha necesaria. Control/acción: Saldos de envases aceptados por cliente y transportista, política V-08 y evidencia de M8. & S1=4; O1=2; D1=2; NPR1=16. D1=2: acta o ensayo específico previo a habilitación; se invalida la mejora de detección si no existe evidencia. & Mantener saldos afectados pendientes; Finanzas resuelve antes de su carga, sin asumir saldo cero. Dueño: DAT; verifica CAL. \\
R-006 & S=4; O=4; D=5; NPR=80. & V-10 no obtiene horas o montos históricos verificables. Dado que las comparaciones históricas de aceptación dependen de horas verificables de entrega y montos esperados de rendición que la captura futura no reconstruye, entonces puede que falten horas verificables de entrega o montos esperados de rendición en los antecedentes históricos, teniendo como consecuencia que no puedan respaldarse con esos registros las comparaciones de OTIF o cuadratura y deba repetirse la validación del período de referencia. & D0=5: faltante de origen o interfaz no verificable puede permanecer oculto hasta conciliar/aceptar. Control/acción: Fuentes históricas de horas y montos verificadas con Comercial, Operaciones y Finanzas, V-10. & S1=4; O1=3; D1=3; NPR1=36. D1=3: aun conciliando, la cobertura histórica puede ser parcial; no se reconstruye lo inexistente. & Plan adicional: JP reúne áreas dentro de dos días hábiles desde bloqueo; en cinco días hábiles exige evidencia o eleva al Comité Ejecutivo la condición contractual. No sustituir historia por mediciones futuras ni aceptar resultado sin referencia. Dueño: IMP; verifica CAL. \\
R-009 & S=4; O=3; D=4; NPR=48. & Calidad no aprueba parámetros por producto al preparar certificación. Dado que la certificación de la parametrización térmica requiere una matriz de umbrales y duraciones por producto aprobada por Calidad del cliente, entonces puede que Calidad no apruebe esa matriz antes del cierre de las pruebas y la marcha blanca de los productos afectados, teniendo como consecuencia que no pueda completarse la evidencia de parametrización y aceptación del alcance térmico y deban reprogramarse validaciones. & D0=4: la validación distribuida puede llegar después de la fecha necesaria. Control/acción: Matriz térmica V-11 versionada y firmada por Calidad por familia antes de certificación. & S1=4; O1=2; D1=2; NPR1=16. D1=2: acta o ensayo específico previo a habilitación; se invalida la mejora de detección si no existe evidencia. & Sin parámetros aprobados, bloquear certificación del alcance afectado y escalar a D-5; no relajar bloqueo térmico. Dueño: CAL; verifica CAL. \\
R-011 & S=4; O=3; D=4; NPR=48. & Acuerdo operativo de un transportista pendiente antes de su ola. Dado que la incorporación de cada empresa transportista al reparto requiere que su acuerdo operativo esté firmado antes de la ola correspondiente, entonces puede que una empresa transportista no firme el acuerdo a tiempo, teniendo como consecuencia que sus conductores no puedan incorporarse ni completar las pruebas de esa ola y quede pendiente demostrar la cobertura del transportista. & D0=4: la validación distribuida puede llegar después de la fecha necesaria. Control/acción: Acuerdo operativo firmado y nómina de cada transportista antes de incorporar la ola. & S1=4; O1=2; D1=2; NPR1=16. D1=2: acta o ensayo específico previo a habilitación; se invalida la mejora de detección si no existe evidencia. & A D-5 sin firma, escalar por CT el mismo día; no habilitar empresa pendiente ni reducir cobertura obligatoria. Dueño: IMP; verifica CAL. \\
R-013 & S=4; O=4; D=4; NPR=64. & Transacciones se siguen registrando solo por el procedimiento anterior. Dado que la implantación exige que los usuarios incorporen sus transacciones al registro oficial de la solución, entonces puede que parte de los usuarios continúe registrando exclusivamente en el procedimiento anterior durante una ola, teniendo como consecuencia que no pueda acreditarse la cobertura de adopción requerida para cerrar esa ola y deba extenderse la intervención de implantación. & D0=4: la validación distribuida puede llegar después de la fecha necesaria. Control/acción: Transacciones de muestra registradas en el flujo nuevo y cobertura certificada por turno/ola. & S1=4; O1=2; D1=2; NPR1=16. D1=2: acta o ensayo específico previo a habilitación; se invalida la mejora de detección si no existe evidencia. & Si persiste flujo paralelo, IMP abre refuerzo en el siguiente turno y mide diariamente; sin evidencia de adopción no cerrar ola. Dueño: IMP; verifica CAL. \\
R-016 & S=4; O=3; D=4; NPR=48. & Decisión requerida no llega en la fecha acordada y bloquea la tarea. Dado que los diseños, pruebas y aceptaciones incluyen actividades que dependen de validaciones o decisiones de áreas del cliente, entonces puede que el área requerida no entregue una validación o decisión en la fecha acordada, teniendo como consecuencia que se detengan o repitan tareas dependientes, se consuma holgura y deban reprogramarse recursos. & D0=4: la validación distribuida puede llegar después de la fecha necesaria. Control/acción: Matriz de decisión con titular y suplente; recepción de insumos y resolución trazables. & S1=4; O1=2; D1=2; NPR1=16. D1=2: acta o ensayo específico previo a habilitación; se invalida la mejora de detección si no existe evidencia. & D-5 sin decisión: escalar un día hábil después; promoción regional: suplente a 15 min y decisión dentro de 30 min, sin actuar por silencio. Dueño: JP; verifica CAL. \\
R-028 & S=4; O=3; D=4; NPR=48. & Acuerdo de finalidades y controles pendiente a D-5 de la ola, u objeción formal al uso operacional propuesto.. Dado que existe objeción sindical a cámaras en cabina y al control de jornada por GPS, aunque la propuesta excluye cámaras en cabina y separa la telemetría operacional, entonces puede que no se acuerden o comprendan las finalidades y controles del uso propuesto antes de la ola de reparto, teniendo como consecuencia resistencia a la incorporación y pruebas o aceptaciones pendientes. & D0=4: la validación distribuida puede llegar después de la fecha necesaria. Control/acción: Acuerdo 5.4.2 con finalidades, accesos y controles comprendidos por Operaciones/TI/sindicato. & S1=4; O1=2; D1=2; NPR1=16. D1=2: acta o ensayo específico previo a habilitación; se invalida la mejora de detección si no existe evidencia. & D-5 sin acuerdo: JP escala el mismo día y conserva sólo uso operacional autorizado; no cámaras en cabina ni GPS de jornada no acordado. Dueño: JP; verifica CAL. \\
R-029 & S=4; O=4; D=4; NPR=64. & Preparador nuevo asignado sin tutor o evaluación, o aprendizaje que excede dos horas en el programa 7.1.2.. Dado que preparación nocturna tiene rotación anual del 38\% y requiere capacitación continua, entonces puede que las incorporaciones superen la capacidad de tutoría y evaluación o lleguen al flujo sin competencias verificadas, teniendo como consecuencia errores de preparación, refuerzo extraordinario y certificación de la ola pendiente. & D0=4: la validación distribuida puede llegar después de la fecha necesaria. Control/acción: Altas por turno, tutor, evaluación antes de autonomía y aprendizaje dentro de dos horas según 7.1.2. & S1=4; O1=3; D1=2; NPR1=24. D1=2: acta o ensayo específico previo a habilitación; se invalida la mejora de detección si no existe evidencia. & Plan adicional: Bodega asigna cobertura certificada antes del turno; IMP/SRE revisa cada alta diariamente. Si supera tutoría, no asignar trabajo autónomo; presentar refuerzo en dos días hábiles y revalidar antes de habilitar. Dueño: IMP; verifica CAL. \\
R-030 & S=4; O=3; D=5; NPR=60. & No hay contrato de interfaz o extracción verificado al cierre del mes 1, o el ensayo temprano contradice formato, acceso o comportamiento declarados.. Dado que el CLIENTE no tiene documentación técnica de las interfaces del ERP ni del WMS de 2013, que se reemplaza progresivamente, entonces puede que el levantamiento no permita verificar a tiempo contratos de intercambio y extracción del legado, teniendo como consecuencia reproceso de adaptadores o cargas y bloqueo de pruebas de integración o migración. & D0=5: faltante de origen o interfaz no verificable puede permanecer oculto hasta conciliar/aceptar. Control/acción: Contrato de extracción e intercambio confirmado con TI; ensayo temprano sin modificar ERP y sustitución progresiva de WMS. & S1=4; O1=3; D1=2; NPR1=24. D1=2: acta o ensayo específico previo a habilitación; se invalida la mejora de detección si no existe evidencia. & Plan adicional: ARQ eleva brecha al finalizar semana 2 del mes 1; TI entrega acceso o alternativa verificable antes de integración. Ensayar archivo controlado si satisface requisitos; bloquear certificación si persiste interfaz no verificable. Dueño: ARQ; verifica CAL. \\
R-031 & S=4; O=3; D=4; NPR=48. & Incidente alto/crítico, evidencia o certificación pendiente en las semanas finales del peak, sin ventana de intervención o validación disponible.. Dado que el peak de septiembre casi duplica entregas durante tres semanas y coincide con el cierre de la marcha blanca E2 en el calendario supuesto de S7, entonces puede que un incidente, una validación pendiente o la falta de disponibilidad de usuarios no pueda resolverse dentro de las ventanas autorizadas, teniendo como consecuencia incumplimiento de condiciones de cierre y necesidad de recuperación de la programación. & D0=4: la validación distribuida puede llegar después de la fecha necesaria. Control/acción: Disponibilidad y continuidad acordadas antes de septiembre; seguimiento diario de las seis condiciones de cierre. & S1=4; O1=3; D1=2; NPR1=24. D1=2: acta o ensayo específico previo a habilitación; se invalida la mejora de detección si no existe evidencia. & Plan adicional: cerrar habilitaciones en mes 19 antes del congelamiento; JP revisa incidentes diariamente en 19--20. Escalar el mismo día una condición incumplida y acordar recuperación en dos días hábiles, sin intervención prohibida ni aceptación ficticia. Dueño: JP; verifica CAL. \\
R-032 & S=4; O=3; D=4; NPR=48. & Aviso de soporte/compatibilidad que deja una versión necesaria fuera de la política aprobada o hace fallar la regresión de la actualización prevista.. Dado que la propuesta opera durante 56 meses con componentes como Laravel 13, PHP 8.5 y Proxmox, entonces puede que cambien su soporte o compatibilidad de forma que las actualizaciones ordinarias no basten, teniendo como consecuencia una adecuación extraordinaria y revalidación para mantener seguridad y continuidad. & D0=4: la validación distribuida puede llegar después de la fecha necesaria. Control/acción: Inventario de versiones y soporte, revisión trimestral y compatibilidad probada antes del límite verificado. & S1=4; O1=2; D1=2; NPR1=16. D1=2: acta o ensayo específico previo a habilitación; se invalida la mejora de detección si no existe evidencia. & Cambio excepcional: ARQ registra aviso en un día hábil y presenta transición en cinco; programar antes de fin de soporte, conservando ventanas y pruebas. Dueño: ARQ; verifica CAL. \\
R-035 & S=4; O=3; D=4; NPR=48. & Conductor asignado y habilitado rechaza el flujo digital o entrega sin la evidencia requerida, sin que la causa sea firma pendiente o reemplazo sin aviso.. Dado que los conductores de terceros no son trabajadores de Puelche y requieren incorporación acordada con su transportista, entonces puede que un conductor autorizado y convocado rechace utilizar el terminal o el flujo digital aun existiendo acuerdo de la empresa, teniendo como consecuencia entregas sin evidencia digital completa y cobertura de adopción de la ola pendiente. & D0=4: la validación distribuida puede llegar después de la fecha necesaria. Control/acción: Demostración práctica en andén e incorporación acordada con transportista y conductor habilitado. & S1=4; O1=2; D1=2; NPR1=16. D1=2: acta o ensayo específico previo a habilitación; se invalida la mejora de detección si no existe evidencia. & Ante rechazo actuar antes de salida: IMP solicita conductor certificado/dispuesto ese turno y registra cobertura pendiente; refuerzo antes de siguiente asignación. Dueño: IMP; verifica CAL. \\
\end{tablalafrox}
\end{landscape}
\endgroup
